% Section module: one-variable encoding of finite geometric fibers.
\section{One witness for a constructible finite-fiber projection}
\label{sec:one-finite-variable}

A constructible set is a finite union of locally closed algebraic sets.
The projection has finite geometric fibers if every fiber over an
algebraic closure has only finitely many points. Finite rational fibers
alone do not suffice for the argument. The construction below needs only
one extra coordinate, including the conditions that denominators be
nonzero; no useful degree or coefficient-height bound is claimed.

\begin{theorem}\label{thm:finite-fibre-one-variable}
Let $X\subseteq\mathbb A^{d+m}_{\mathbb Q}$ be constructible over
$\mathbb Q$, and suppose the coordinate projection
$\pi:X\to\mathbb A^d_{\mathbb Q}$ has finite geometric fibers.
There is an effectively constructible polynomial
$P\in\mathbb Z[\mathbf x,U]$ such that, for every $\mathbf a\in\mathbb Q^d$,
\[
 \mathbf a\in\pi(X(\mathbb Q))
 \quad\Longleftrightarrow\quad
 \exists U\in\mathbb Q:\ P(\mathbf a,U)=0.
\]
\end{theorem}
\begin{proof}
We give the construction including its exceptional loci. Decompose $X$
into finitely many locally closed affine pieces; represent an inequation
$h\ne0$ by the coordinate relation $hv=1$ when forming a coordinate ring.
The added coordinate is uniquely determined, so geometric fibers stay
finite. Decompose the image closure into irreducible pieces and restrict
accordingly. On one piece let $R$ be the coordinate ring of the irreducible
base closure $B$, let $L=\operatorname{Frac}(R)$, and let $A$ be the reduced
coordinate ring of the affine source. Its generic fiber algebra
$A_L$ is finite-dimensional, reduced, and separable over $L$; it is a
finite product of finite separable field extensions.

The primitive element theorem for a finite separable algebra gives
$A_L\simeq L[T]/(h)$ with $h$ monic. Clearing the denominators of a
primitive element, of expressions for the source generators in that
primitive element, and of their finitely many relations gives a nonzero
$g\in R$ for which
\[
 A_g\simeq R_g[T]/(h),\qquad h\in R_g[T]\text{ monic}.
\]
For precision, the map is surjective after clearing the generator
expressions; its kernel vanishes over $L$ and is therefore zero because
$R_g[T]/(h)$ is free over the domain $R_g$. Any intermediate relation
which vanishes only over $L$ is made zero by a further localization.
Thus this is an actual isomorphism on an open base set, not just a
bijection on generic points. Consequently a rational point of $B\cap D(g)$
has a rational lift exactly when the specialized monic $h$ has a rational
root. The residual part lies over a proper closed subset of $B$, so
induction on the dimension of the image closure gives finitely many such
charts. Each chart comes with polynomial equations $E_i(\mathbf x)=0$
for its base closure and one or more polynomial inequations; multiply the
latter together to obtain $G(\mathbf x)\ne0$. Choose polynomial
representatives in the ambient base space for all elements of $R$.

It remains to express each chart with one quantified variable. Write its
monic equation, after choosing a common coefficient denominator, as
\[
 h(T)=T^n+\sum_{j=1}^n\frac{N_j(\mathbf x)}{D(\mathbf x)}T^{n-j}.
\]
Add $D\ne0$ to the recorded chart conditions. A rational root of $h$
may be zero; to cover this exactly, use the two fixed alternatives
$h(T)$ and $h(T-1)$. If $r$ is a rational root of $h$, then the two
polynomials have roots $r$ and $r+1$, at least one of which is nonzero.
Conversely a rational root of either polynomial gives a rational root of
$h$. The shifted polynomial is still monic and has the same permissible
common denominator $D$.

For either alternative, with its coefficient numerators again denoted
$N_j$, set
\[
 \Psi(\mathbf x,U)=1+\sum_{j=1}^n
 N_j(\mathbf x)D(\mathbf x)^{2j-1}G(\mathbf x)^jU^j.
\]
At $DG=0$ this polynomial is identically $1$. At $DG\ne0$, put
$L_0=D^2G$. Then, for $U\ne0$,
\[
 \Psi(\mathbf x,U)=(L_0U)^n h(1/(L_0U)),
\]
where $h$ now denotes the selected alternative. Also
$\Psi(\mathbf x,0)=1$. Hence a rational zero of $\Psi$ exists exactly
when all inequations hold and that alternative has a nonzero rational
root. This construction cannot acquire spurious solutions at a vanished
denominator, because its value there is exactly one.

The polynomial
\[
 Q(\mathbf x,U)=\Psi(\mathbf x,U)^2+\sum_iE_i(\mathbf x)^2
\]
therefore describes that chart and alternative. Multiply $Q$ over the
finite list of charts and the two alternatives, always reusing the same
variable $U$, and clear numerical coefficient denominators. This gives
$P$. A product is zero if and only if some factor is zero, and over
$\mathbb Q$ each displayed sum of squares is zero if and only if all of
its summands vanish. The empty image is represented by $P=1$.
All the steps are effective operations in polynomial rings over
$\mathbb Q$, including ideal decomposition and localization, followed by
finite-dimensional algebra over the resulting function fields. The
induction terminates because the dimension of every exceptional base
strictly decreases. This effectivity assertion does not mean that these
computations have been executed for the source system.
\end{proof}

\begin{corollary}
A finite-fiber presentation retaining $d$ rational coordinates, one of
which is the designated input, yields a polynomial in at most $d+1$
total variables, with at most $d$ existential witnesses and the same
rational input projection. In particular, the presentation with $d=557$
proved below gives at most $558$ total variables. Its identification
with $\mathbb Z$ retains the arithmetic hypothesis of the source.
\end{corollary}

The elementary reciprocal transformation in this proof agrees with
\cite[Section 10.6, equation (10.8)]{Sun2026Rational86}; the chart argument
above extends it from the finite algebras explicitly treated there to all
constructible sets with finite geometric fibers. This citation supports
the algebraic device only. The bounded-variable undecidability theorem
in that recent preprint and its analytic arithmetic inputs are not used.
Its convention is that a monic degree-$n$ polynomial is written
$T^n+\sum_{j=1}^n(N_j/D)T^{n-j}$: the reciprocal polynomial has exponent
$U^j$ and denominator-clearing factor $D^{2j-1}$, not $D^{j-1}$.
The two shifts are essential if zero could be the only rational root.
All mathematical claims used here are proved above; no result about
the arithmetic of that preprint is needed.
